% !TeX Program = lualatex

\documentclass[a4paper]{article}
\usepackage[left=2cm,right=2cm,top=1cm,bottom=1cm]{geometry}
\usepackage{fontspec}
\usepackage{libertinus}
\renewcommand*\familydefault{\sfdefault}

\usepackage{xcolor}
\colorlet{textblue}{black!40!blue!50!white}
\colorlet{bgblue}{black!40!blue!10!white}
\colorlet{xblue}{black!40!blue!25!white}
\colorlet{lineblue}{black!40!blue}

\usepackage{pst-barcode}
\usepackage{pifont}

\usepackage{tikz}
\usetikzlibrary{calc}
\usetikzlibrary{shapes.geometric}
\usetikzlibrary{decorations.pathmorphing}

\usepackage{enumitem}

\usepackage{tabularx}
\usepackage{colortbl}
\usepackage{booktabs}
\usepackage{array}
\arrayrulecolor{lineblue}
\setlength{\aboverulesep}{0pt}
\setlength{\belowrulesep}{0pt}
\newcommand{\PreserveBackslash}[1]{\let\temp=\\#1\let\\=\temp}
\newcolumntype{C}[1]{>{\PreserveBackslash\centering}p{#1}}
\newcolumntype{R}[1]{>{\PreserveBackslash\raggedleft}p{#1}}
\newcolumntype{L}[1]{>{\PreserveBackslash\raggedright}p{#1}}
\newcolumntype{Y}{>{\centering\arraybackslash}X}

\usepackage{calc}
\usepackage{everypage}
\usepackage{lastpage}
\usepackage{multicol}

\def\documentid{F-BLAU-2k26}

\AddEverypageHook{%
    \begin{tikzpicture}[remember picture, overlay]
        \node [color=textblue, rotate=90, anchor=north] at ($(current page.west)+(10mm,0)$)
        {\ttfamily\bfseries Originalpapier nur, wenn dieser Hinweis in Blaudruck erscheint.};
        \node [rotate=90, anchor=north west, align=left] at ($(current page.south west)+(10mm,10mm)$)
        {%
            \hspace{1cm}\texttt{\documentid}
            \\
            \hspace{1cm}\texttt{Seite \thepage{} von \pageref{LastPage}}
            \\
            \psbarcode[transy=0.35,scalex=0.5,scaley=0.5]{\documentid-\thepage}{parse,format=square: version=18x18}{datamatrix}
        };
    \end{tikzpicture}%
    %
    \def\foldFirst{105mm}%
    \def\foldSecond{210mm}%
    \begin{tikzpicture}[remember picture, overlay]
        \node (foldAA) at ($(current page.north west) + (5mm, - \foldFirst)$) {};
        \node (foldAB) at ($(current page.north west) + (10mm, - \foldFirst)$) {};
        \node (foldBA) at ($(current page.north west) + (5mm, - \foldSecond)$) {};
        \node (foldBB) at ($(current page.north west) + (10mm, - \foldSecond)$) {};
        \node (centerA) at ($(current page.west) + (5mm, 0)$) {};
        \node (centerB) at ($(current page.west) + (10mm, 0)$) {};
        \draw[-] (foldAA) -- (foldAB);
        \draw[-] (foldBA) -- (foldBB);
        \draw[-] (centerA) -- (centerB);
    \end{tikzpicture}%
}

\parindent=0pt
\pagestyle{empty}

\usepackage[ngerman]{babel}

\newcommand\sectitle[1]{\textbf{\Large\raisebox{0pt}[1.5\height][1.5\depth]{\vphantom{#1}}#1}}
\newcommand\tftitle[1]{\textbf{#1}}
\newcommand\tfnote[1]{\textit{#1}}
\newcommand\X{\begin{tabular}{|>{\columncolor{white}}l|}\hline\raisebox{0pt}[1.5\height][1.5\depth]{\vphantom{Fg}\hspace{-0.29em}\color{xblue}X\hspace{-0.29em}}\\\hline\end{tabular}}
\newcommand\digittf{\begin{tabular}{|>{\columncolor{white}}l|}\hline\raisebox{0pt}[1.5\height][1.5\depth]{\vphantom{Fg}}\\\hline\end{tabular}}
\newcommand\widetf{\begin{tabular}{|>{\columncolor{white}}l|}\hline\raisebox{0pt}[1.5\height][1.5\depth]{\vphantom{Fg}}\hspace*{\linewidth - 1.5em}\\\hline\end{tabular}}
\newcommand\Xsub[1]{$\underset{#1}{\X}$}

\newcounter{ansnum}
\newcommand\ans{\refstepcounter{ansnum}\arabic{ansnum}}

\begin{document}%
\begin{tikzpicture}[remember picture, overlay]
    \node [rotate=90, anchor=south] at ($(current page.east)+(-15mm,6.75cm)$)
    {\ttfamily\bfseries\footnotesize Dieses Formular ist auch auf grünem Papier blau.};
    \node [anchor=south west] at ($(current page.south west)+(0,5mm)$)
    {%
        \hspace*{1cm + 0.33ex}%
        \ttfamily%
        CC BY-SA 4.0
        \cdot{}
        Merkblatt der Hafenkommandantur umseitig
    };
\end{tikzpicture}%
Blaues Formular zur Beantragung des Antragsformulars für den

\begin{Huge}
    Passierschein A38
\end{Huge}

~

\begin{tabularx}{\linewidth}{@{}Xr@{}}
    \tftitle{Laufende Nummer} \tfnote{(Laufende Nummer des A38 rückwärts)}                           & Ausgabestelle: Schalter I         \\
    \digittf\digittf\digittf\digittf\digittf\digittf\digittf\digittf\digittf\digittf\digittf\digittf & nach Rundschreiben B65
\end{tabularx}

\vspace*{3ex}

\hspace*{-1cm}\begin{tabularx}{\linewidth + 1cm}{@{} R{7.5mm} @{\hspace*{2.5mm}}| >{\columncolor{bgblue}}X |@{}}
    \cmidrule{2-2}
         & \sectitle{Angaben zum Antrag auf den Antrag}                                                  \\
    \ans & \tftitle{Nennen Sie Ihre Lieblingsfarbe.}                                                     \\
         & \tfnote{Andere Antworten als „Blau“ werden an Schalter XII, Stiege B, Korridor J verwiesen.}  \\
         & \widetf                                                                                       \\
         &                                                                                               \\
    \ans & \tftitle{Wie viele Schalter haben Sie heute bereits aufgesucht?}                             \\
         & \tfnote{Bitte in römischen Ziffern. Schalter I zählt nicht, da Sie gerade dort stehen.}       \\
         & \digittf\digittf\digittf\digittf\digittf\digittf                                              \\
         &                                                                                               \\
    \ans & \tftitle{Wie lange warten Sie bereits?}                                                       \\
         & \X{} Minuten \qquad \X{} Stunden \qquad \X{} Mondphasen \qquad \X{} Ich bin hier aufgewachsen. \\
         &                                                                                               \\
    \ans & \tftitle{Wie bereitwillig steigen Sie Treppen?}                                               \\
         & Ungern \Xsub{1} \Xsub{2} \Xsub{3} \Xsub{4} \Xsub{5} Bis in den sechsten Stock                \\
         &                                                                                               \\
    \cmidrule{2-2}
         & \sectitle{Angaben für die Hafenkommandantur}                                                  \\
    \ans & \tftitle{Wie gedenken Sie, den Hafen zu erreichen?}                                           \\
         & \X{} Galeere \qquad \X{} Floß \qquad \X{} Fass \qquad \X{} Schwimmend \tfnote{(Anlage N beifügen)} \\
         &                                                                                               \\
    \ans & \tftitle{Führen Sie Ladung mit, die schwerer ist als Ihr Wasserfahrzeug?}                     \\
         & \X{} Nein \qquad \X{} Ja, einen Stein \qquad \X{} Ja, mehrere Steine                          \\
         &                                                                                               \\
    \ans & \tftitle{Gewünschter Liegeplatz:}                                                             \\
         & \tfnote{Liegeplätze werden in der Reihenfolge des Ablegens vergeben.}                        \\
         & \digittf\digittf\digittf                                                                      \\
         &                                                                                               \\
    \ans & \tftitle{Tragen Sie den Wasserstand zum Zeitpunkt des Ausfüllens ein.}                        \\
         & \tfnote{Bei Ebbe bitte zwei Zeilen tiefer eintragen.}                                         \\
         & \widetf                                                                                       \\
         &                                                                                               \\
    \ans & \tftitle{Ist Ihr Wasserfahrzeug in den letzten 14 Monaten:}                                   \\
         & \tfnote{Nur Einfachantwort möglich.}                                                          \\
         & \X{} gesunken                                                                                 \\
         & \X{} versenkt worden                                                                          \\
         & \X{} beides, aber nicht in dieser Reihenfolge                                                 \\
         & \X{} ohne Passierschein A38 in See gestochen                                                  \\
         &                                                                                               \\
    \cmidrule{2-2}
\end{tabularx}

\begin{multicols}{3}
    \tftitle{Hinweise}\\
    \begin{footnotesize}
        Dieses Formular ist nur zusammen mit dem Passierschein A38 gültig. Der Passierschein A38 wird nur gegen Vorlage dieses Formulars ausgestellt.

        Das Formular gilt als blau, solange es nicht nass wird.

        Bei Farbfehlsichtigkeit wenden Sie sich bitte an Schalter Grün.
    \end{footnotesize}

    \vspace*{2.5ex}\tftitle{Stempel des Hafens:}\par
    \begin{tabularx}{\linewidth}{|Y|}
        \midrule
        \\
        \tfnote{(\ans{} Nur bei Flut stempeln.)} \\
        \\
        \midrule
    \end{tabularx}

    \columnbreak

    \tftitle{Bitte fügen Sie folgende Anlagen bei:}
    \begin{itemize}[noitemsep,topsep=0pt,parsep=0pt,partopsep=0pt]
        \item Passierschein A38 \tfnote{(siehe dort)}.
        \item Bescheinigung über das Nichtvorliegen des rosa Formulars.
        \item Liegeplatzquittung der Hafenkommandantur \tfnote{(umseitig)}.
        \item Bestätigung des Pförtners, dass die antragstellende Person das Gebäude betreten hat
        \item Wartemarke des Erdgeschosses, ungeknickt
    \end{itemize}

    \begin{tabularx}{\linewidth}{@{}lX@{}}
        \ans{} \tftitle{Daumenabdruck}
         &
        \tikz \draw[dashed,lineblue] (0,0.6) ellipse (0.45 and 0.65);
    \end{tabularx}

    \columnbreak

    \begin{itemize}[noitemsep]
        \item[\ans{} \X] Ich bestätige, dass dieses Formular blau ist.
        \item[\ans{} \X] Ich habe Rundschreiben B65 nicht gelesen. \tfnote{(Pflichtfeld)}
        \item[\ans{} \X] Seemannsgarn?
            \\[4ex]
        \item[\ans{} \phantom{\X}] \dotfill\\
            Unterschrift (blau)
            \smash{
                \begin{tikzpicture}[remember picture, overlay]
                    \node [
                        draw,
                        dashed,
                        lineblue,
                        ellipse,
                        align=center,
                        minimum width=6cm,
                        minimum height=3cm,
                        inner sep=0pt,
                        rotate=90
                    ]
                    at ($(current page.east)+(5mm,-11cm)$) {
                        \color{black} \ans{} \\
                        \color{black} Hier zum Schiffchen falten.
                        \\
                        \\
                        \\
                        \\
                        \\
                        \\
                    };
                \end{tikzpicture}
            }
    \end{itemize}%
\end{multicols}%
\clearpage%

\begin{tikzpicture}[remember picture, overlay]
    \node [anchor=south west] at ($(current page.south west)+(0,5mm)$)
    {%
        \hspace*{1cm + 0.33ex}%
        \ttfamily%
        Hafenkommandantur
        \cdot{}
        Geöffnet bei Flut
        \cdot{}
        Bei Ebbe bitte unten klingeln
    };
\end{tikzpicture}%
Merkblatt der

\begin{Huge}
    Hafenkommandantur
\end{Huge}

~

\begin{tabularx}{\linewidth}{@{}Xr@{}}
    \tftitle{Für Antragstellende des Passierscheins A38,} & zuständig für alles,  \\
    \tftitle{die den Weg zum Hafen suchen.}               & was nass werden kann
\end{tabularx}

\vspace*{3ex}

\hspace*{-1cm}\begin{tabularx}{\linewidth + 1cm}{@{} R{7.5mm} @{\hspace*{2.5mm}}| >{\columncolor{bgblue}}X |@{}}
    \cmidrule{2-2}
     & \sectitle{Lageplan}                                                                    \\
     & \tfnote{Nicht maßstabsgetreu. Wasser nicht wassergetreu.}                              \\
     & \begin{tikzpicture}[x=1cm,y=1cm, font=\footnotesize]
           % Stadt
           \fill[white] (0,0) rectangle (15.4,6.4);
           \draw[lineblue] (0,0) rectangle (15.4,6.4);
           \draw[lineblue, fill=bgblue] (0.6,3.6) rectangle (4.4,5.8);
           \node[align=center] at (2.5,4.7) {Haus, das\\Verrückte macht};
           \node[anchor=north west, align=left] at (0.6,3.5) {\tfnote{Sie sind vermutlich}\\\tfnote{noch hier.}};
           % Schalter
           \foreach \i/\x in {I/5.4,XII/6.6} {
               \draw[lineblue, fill=white] (\x,4.9) rectangle ++(0.9,0.6);
               \node at (\x+0.45,5.2) {\i};
           }
           % Wasser
           \fill[bgblue] (9.4,0) .. controls (9.0,2) and (10.4,3.5) .. (9.8,6.4) -- (15.4,6.4) -- (15.4,0) -- cycle;
           \foreach \y in {0.8,1.8,...,5.8} {
               \draw[lineblue, decorate, decoration={snake, amplitude=0.6mm, segment length=4mm}] (11.2,\y) -- (14.8,\y);
           }
           \draw[lineblue] (0,0) rectangle (15.4,6.4);
           \node[align=center, fill=bgblue] at (13,3.3) {\textbf{Hafen}\\\tfnote{immer da, wo Wasser ist}};
           % Mole und Kommandantur
           \draw[lineblue, fill=white] (9.4,1.1) rectangle (11.4,1.5);
           \draw[lineblue, fill=white] (8.0,0.6) rectangle (9.5,1.9);
           \node[align=center] at (8.75,1.25) {\scriptsize Kom-\\[-2pt]\scriptsize mandantur};
           % Weg
           \draw[thick, dashed, black, ->, rounded corners=3mm] (4.4,4.2) -- (5.85,4.2) -- (5.85,4.85);
           \draw[thick, dashed, black, ->] (6.3,5.2) -- (6.58,5.2);
           \draw[thick, dashed, black, ->, rounded corners=3mm] (7.05,4.9) -- (7.05,3.6) -- (5.2,3.6) -- (5.2,2.4) -- (8.75,2.4) -- (8.75,1.95);
           \node[anchor=south, font=\scriptsize\itshape] at (7.0,2.4) {Bitte nicht abkürzen};
           \node[anchor=west, align=left] at (0.4,1.6) {\tfnote{Empfohlener Weg.}\\\tfnote{Kürzere Wege sind}\\\tfnote{nicht genehmigt.}};
       \end{tikzpicture}                                                                      \\
     &                                                                                        \\
    \cmidrule{2-2}
     & \sectitle{Hafenordnung (Auszug)}                                                       \\
     & \begin{enumerate}[label=\S\,\arabic*,leftmargin=2.2em,itemsep=0.6ex]
           \item Der Hafen befindet sich unterhalb der Stadt. Wer ihn oberhalb der Stadt antrifft, meldet dies unverzüglich der Hafenkommandantur.
           \item Das Anlegen ist nur mit gültigem Passierschein A38 gestattet. Das Ablegen ist nur ohne gültigen Passierschein A38 gestattet, da dieser beim Anlegen einbehalten wird.
           \item Galeeren werden nach Rudern abgerechnet. Floße nach Stämmen. Fässer nach Inhalt.
           \item Wer im Hafenbecken schwimmt, gilt als Wasserfahrzeug und benötigt einen Liegeplatz.
           \item Steine, Hinkelsteine und vergleichbare Ladung sind beim Zoll anzumelden und dürfen nicht zur Abfertigung anderer Wasserfahrzeuge verwendet werden.
           \item Die Hafenkommandantur haftet nicht für Gezeiten.
       \end{enumerate}                                                                        \\
    \cmidrule{2-2}
\end{tabularx}

\vfill

\begin{tikzpicture}
    \draw[dashed, lineblue] (-1cm,0) -- (\linewidth,0);
    \node[fill=white, inner sep=1pt] at (-0.6cm,0) {\ding{34}};
    \node[anchor=south east, font=\footnotesize\itshape] at (\linewidth,0) {Hier abtrennen und dem blauen Formular beifügen.};
\end{tikzpicture}

\vspace*{1ex}

\begin{tabularx}{\linewidth}{@{}Xr@{}}
    \textbf{\large Liegeplatzquittung}                     & \tftitle{Liegeplatz}                \\
    \tfnote{Gültig von Ankunft bis zur nächsten Ebbe.}     & \digittf\digittf\digittf            \\[1ex]
    \tftitle{Name des Wasserfahrzeugs}                     & \tftitle{Anzahl Ruder}              \\
    \begin{tabular}{|>{\columncolor{white}}l|}\hline\raisebox{0pt}[1.5\height][1.5\depth]{\vphantom{Fg}}\hspace*{9cm}\\\hline\end{tabular}
                                                           & \digittf\digittf\digittf
\end{tabularx}

\vspace*{2ex}
\begin{center}
    \tfnote{(Der Platz unter dieser Zeile ist für Spritzwasser vorgesehen.)}
\end{center}
\vspace*{8ex}

\end{document}
